\documentclass[11pt,a4paper]{article}
\usepackage{curso}
\begin{document}
\portada{Ejercicios --- Lección}{7.8}{Complejidad comparada}
\section{Ejercicios}
\begin{enumerate}[label=\textbf{E\arabic*.}]
\item (B) ¿Qué recurso físico consume cada clase de paso?
\item (B) Calcule $\ell_0$ para $n=2048$ bits.
\item (B) ¿Por qué leer la salida de HHL anula su ventaja para el vector completo?
\item (B) Escriba las cuatro preguntas para una afirmación de ventaja.
\item (I) Estime la exponencial de la criba de cuerpos de números para $n=1024$ y compárela con $n=2048$.
\item (I) Si cada nodo de difusión ocupa $1\ \mu\text{m}^2$, ¿qué área necesita una red de $2^{64}$ nodos? Compárela con la superficie terrestre ($5{,}1\times10^{14}$ m$^2$).
\item (A) Haga una tabla de recursos (tiempo, área, energía, precisión, etiqueta de evidencia) para factorizar $N=299$ con los cuatro métodos de la unidad.
\end{enumerate}
\section{Soluciones}
\begin{solucion}
\textbf{E1.} Cuántico: coherencia y corrección de errores; difusión: área/nodos y precisión analógica; digital: tiempo.

\textbf{E2.} $M=2049$: $4\cdot2050\cdot2050\approx1{,}68\times10^7$.

\textbf{E3.} Estimar las $d$ amplitudes exige del orden de $d$ repeticiones o más: se pierde el factor $\log d$.

\textbf{E4.} Paso y costo; crecimiento de área, energía y precisión; salida y costo de lectura; línea base y condiciones.

\textbf{E5.} $n=1024$: $\approx2^{87}$; $n=2048$: $\approx2^{117}$ (sin $o(1)$).

\textbf{E6.} $2^{64}\cdot10^{-12}\ \text{m}^2\approx1{,}8\times10^7\ \text{m}^2$ ($18\ \text{km}^2$); con $2^{100}$ nodos, $\approx1{,}3\times10^{18}$ m$^2$, miles de veces la superficie terrestre.

\textbf{E7.} Para $N=299$ todos son triviales; el valor del ejercicio es escribir explícitamente cada columna y su etiqueta.
\end{solucion}
\end{document}
